\begin{solapartat}
\figura[0.45]{sol_fig1.png}

\punts{0,3} ${}^{212}_{82}Pb \rightarrow {}^{212}_{83}Bi + {}^{\;\,0}_{-1}\beta + {}^{0}_{0}\bar{\nu}_e$\\
Desintegració beta $\beta^-$ (electró)

\punts{0,3} ${}^{212}_{83}Bi \rightarrow {}^{208}_{81}Tl + {}^{4}_{2}\alpha$\\
Desintegració alfa, $\alpha$ (nuclis d'heli)

\punts{0,2}
\[ {}^{212}_{82}Pb\left\{\begin{aligned} &\text{nom: plom}\\ &\text{nombre atòmic: 82}\\ &\text{nombre màssic: 212}\end{aligned}\right. ;\quad
{}^{212}_{83}Bi\left\{\begin{aligned} &\text{nom: bismut}\\ &\text{nombre atòmic: 83}\\ &\text{nombre màssic: 212}\end{aligned}\right. ;\quad
{}^{208}_{81}Tl\left\{\begin{aligned} &\text{nom: tal·li}\\ &\text{nombre atòmic: 81}\\ &\text{nombre màssic: 208}\end{aligned}\right. \]

\punts{0,2} ${}^{\;\,0}_{-1}\beta = {}^{\;\,0}_{-1}e = \text{electró - partícula } \beta$; ${}^{4}_{2}\alpha = \text{partícula } \alpha$; ${}^{0}_{0}\bar{\nu}_e = \text{antineutrí electrònic}$

Descomptarem \punts{0,1} per cada mancança.
\end{solapartat}

\begin{solapartat}
\punts{0,2} Si s'han desintegrat el $10\,\%$, queden el $90\,\%$ dels nuclis
\[ \rightarrow N = 0,9N_0 \Rightarrow \frac{N}{N_0} = 0,9 \]

Obtenim la constat de desintegració radioactiva, $\lambda$, a partir del període de semidesintegració que es el temps que ha de passar per reduir-se a la meitat la quantitat d'una substancia radioactiva

\punts{0,2} $\dfrac{1}{2}N_0 = N_0e^{-\lambda T_{1/2}} \Rightarrow \lambda = \dfrac{\ln 2}{T_{1/2}}$

Calculem el temps que haurà de passar per tal de que es desintegrin un $10\,\%$ dels nuclis:

\punts{0,2} $t = -\left(\dfrac{T_{1/2}}{\ln 2}\right)\left(\ln\dfrac{N}{N_0}\right)$

\punts{0,4} $t = -\left(\dfrac{10,64}{\ln 2}\right)(\ln 0,9) = 1,617\ h$
\end{solapartat}
